%%% Copyright Sheldon Axler, 2024. This file may be used only with the permission of Sheldon Axler.


\chapter{Vector Spaces} \label{VectorSpaces} %\label{129}



Linear algebra is the study of linear maps on finite-dimensional vector spaces.  Eventually we will learn what all these terms mean.  In this chapter we will define vector spaces and discuss their elementary properties.

In linear algebra, better theorems and
more insight emerge if  complex numbers are investigated along with real
numbers.  Thus we will begin by introducing the complex numbers and their basic
properties.

We will generalize the examples of a plane and of ordinary space to $\R{n}$ and $\CC{n}$, which we then will generalize to the notion of a vector space. As we will see, a vector space is a set with operations of addition and scalar multiplication that satisfy natural algebraic properties.

Then our next topic will be subspaces, which play a role for vector spaces analogous to the role played by subsets for sets. Finally, we will look at sums of subspaces (analogous to unions of subsets) and direct sums of subspaces (analogous to unions of disjoint sets).

\begin{comment}
\TopicsBox{
\item
basic properties of the complex numbers
\item
$\R{n}$ and $\C{n}$
\item
vector spaces
\item
subspaces
\item
sums and direct sums of subspaces
}
\end{comment}


\FigureHereCredit{artwork/\FigureName}{\FigureScale}%
{Ren\'e Descartes\index{Descartes, Ren\'e} explaining his work to Queen Christina of Sweden\index{Christina, Queen of Sweden}.\\Vector spaces are a generalization of the description of a plane\\using two coordinates, as published by Descartes in 1637.}{0.8in}{Pierre Louis Dumesnil, Nils Forsberg}

%%\ifnum\sizepaper=2
%%    \section{\texorpdfstring{$\R n$ and $\C n$}{R\unichar{"005E}n and C\unichar{"005E}n}} % To get correct pdf bookmarks.
%%    \else
    \section{$\R n$ and $\C n$}
\label{Rn_and_Cn}
%%    \fi

%\section{$\R{n}$ and $\C{n}$}

\subsection{Complex Numbers}

You should already be familiar with basic properties of the set
$\R{}$\sindex{R@$\R{}$} of real numbers.  Complex numbers were invented so that we
can take square roots of negative numbers.  The
idea is to assume we have  a square root of $-1$,
denoted by $i$, that obeys the usual rules of arithmetic. Here are the formal definitions.

\DefinitionBox{complex numbers, $\symbf{C}{}$}{%
\begin{Itemize}%
\item
A \emph{complex number}\index{complex number} is an ordered pair
$(a, b)$, where $a, b \in \R{}$, but we will write this as $a + bi$.

\item
The set of all complex numbers is denoted by $\C{}$\sindex{C@$\C{}$}:
\[
\C{} = \{ a + bi: a, b \in \R{} \}.
\]
\vspace{-.27in}
\item
\emph{Addition} and \emph{multiplication}\index{addition!of complex numbers}\index{product!of complex numbers} on $\C{}$ are defined by
\begin{gather*}
(a + bi) + (c + di) = (a + c) + (b + d)i, \\
(a + bi)(c + di) = (ac - bd) + (ad + bc)i;
\end{gather*}
here $a, b, c, d \in \R{}$.
\end{Itemize}}

If $a \in \R{}$, we identify $a + 0i$ with the real number $a$.    Thus we
think of $\R{}$ as a subset of $\C{}$. We usually write $0 + bi$ as just $bi$, and we usually write $0 + 1i$ as just~$i$.

\mar{The symbol\/ $i$ was first used to denote\/ $\sqrt{-1}$ by Leonhard Euler in 1777.}

To motivate the definition of complex multiplication given above, pretend that we knew that $i^2 = -1$ and then use the usual rules of arithmetic to derive the formula above for the product of two complex numbers. Then use that formula to verify that we indeed have
\[
i^2 = -1.
\]

Do not memorize the formula for the product of two complex numbers---you can always rederive it by recalling that
$i^2 = -1$  and then using the usual rules of arithmetic (as given by \ref{complexarith}). The next example illustrates this procedure.

\begin{example}{complex arithmetic}
The product $(2 + 3i)(4+5i)$ can be evaluated by applying the distributive and commutative properties from \ref{complexarith}:
\begin{align*}
(2 + 3i)(4+5i) &= 2 \cdot (4+5i) +(3i)(4+5i) \\
&= 2 \cdot 4 + 2 \cdot 5i + 3i\cdot 4 + (3i)(5i) \\
&= 8 + 10i + 12i -15 \\
&= -7 + 22i.\\[-24bp]
\end{align*}
\end{example}

Our first result states that complex addition and complex multiplication have the familiar properties that we expect.

\ResultBox[complexarith]{properties of complex arithmetic}{%
\begin{description}%
\item[commutativity] \ \\
$\al + \beta = \beta + \al$ and $\al \beta = \beta \al$ for all $\al, \beta \in \C{}$.\index{commutativity}

\item[associativity] \ \\
$(\al + \beta) +  \la = \al + (\beta + \la)$ and
$(\al \beta)\la = \al(\beta \la)$ for all $\al, \beta, \la \in \C{}$.\index{associativity}

\item[identities] \ \\
$\la + 0 = \la$ and $\la 1 = \la$ for all $\la \in \C{}$.

\item[additive inverse] \ \\
For every $\al \in \C{}$, there exists a unique $\beta \in \C{}$ such that
$\al + \beta = 0$\index{additive inverse!in $\C{}$}.

\item[multiplicative inverse] \ \\
For every $\al \in \C{}$ with $\al \ne 0$, there exists a unique $\beta \in \C{}$ such that $\al \beta = 1$.

\item[distributive property] \ \\
$\la(\al + \beta) = \la \al + \la \beta$ for all $\la, \al, \beta \in \C{}$.\index{distributive property}
\end{description}}

The properties above are proved using the familiar properties of real numbers and the definitions of complex addition and multiplication. The next example shows how commutativity of complex multiplication is proved. Proofs of the other properties above are left as exercises.

\begin{example}{commutativity of complex multiplication}\index{commutativity}
To show that $\al \beta = \beta \al$ for all $\al, \beta \in \C{}$,
suppose
\[
\al = a + bi \andd \beta = c + di,
\]
where $a, b, c, d \in \R{}$. Then the definition of multiplication of complex numbers shows that
\begin{align*}
\al \beta &= (a + bi)(c + di) \\
&= (ac - bd) + (ad + bc)i
\end{align*}
and
\begin{align*}
\beta\al &= (c + di)(a + bi) \\
&= (ca - db) + (cb + da)i.
\end{align*}
The equations above and the commutativity of multiplication and addition of real numbers show that $\al \beta = \beta \al$.
\end{example}

Next, we define the additive and multiplicative inverses of complex numbers, and then use those inverses to define subtraction and division operations with complex numbers.

\DefinitionBox{$-\al$, subtraction, $1/\al$, division}{
Suppose $\al, \beta \in \C{}$.
\begin{Itemize}
\item
Let $-\al$ denote the additive inverse\index{additive inverse!in $\C{}$} of $\al$.  Thus $-\al$ is the
unique complex number such that
\[
\al + (-\al) = 0.
\]

\item
\emph{Subtraction}\index{subtraction of complex numbers} on $\C{}$ is defined by
\[
\beta - \al = \beta + (-\al).
\]

\item
For $\al \ne 0$, let $1/\al$ and $\frac{1}{\al}$ denote the multiplicative inverse
of~$\al$.  Thus $1/\al$ is the unique complex number such that
\[
\al (1/\al) = 1.
\]
\item
For $\al \ne 0$, \emph{division} by $\al$ is defined by\index{division of complex numbers}
\[
\beta / \al = \beta (1/\al).
\]
\end{Itemize}}

So that we can conveniently make definitions and prove theorems that apply to both real and complex numbers, we adopt the following notation.

%\NotationBox{$\F{}$}{Throughout this book,
\NotationBox{}{Throughout this book,
$\F{}$ stands for either $\R{}$ or $\C{}$.}
\sindex{F@$\F{}$}

\mar[4]{The letter\/ $\F{}$ is used because\/
$\R{}$ and\/ $\C{}$ are examples of what  are called\/ \marbf{fields}.}

Thus if we prove a theorem involving $\F{}\3$, we will know that it
holds when $\F{}$ is replaced with $\R{}$ and when $\F{}$ is replaced with $\C{}$.

Elements of $\F{}$ are called \emph{scalars}\index{scalar}.  The word
``scalar'' (which is just a fancy word for ``number'') is often used when we want to emphasize that
an object is a number, as opposed to a vector (vectors will be defined soon).

For $\al \in \F{}$ and $m$ a positive integer, we define $\al^m$ to denote the
product of $\al$ with itself $m$ times:
\[
\al^m = \underbrace{\al \cdots \al}_{m \text{ times}}\!\!.
\]
This definition implies that
\[
\paren[\big]{\al^m}^n = \al^{mn} \andd (\al \beta)^m = \al^m \beta^m
\]
for all $\al, \beta \in \F{}$ and all positive integers $m, n$.

\pagebreak[3]

\subsection{Lists}

Before defining $\R{n}$ and $\CC{n}$, we look at two important examples.

\begin{example}{$\R{2}$ and\/ $\R{3}$}
\begin{Itemize}
\item
The set $\RR{2}$, which you can think of as a plane, is the set of all ordered pairs of real numbers:
\[
\R{2} = \br[\big]{(x,y): x, y \in \R{} }.
\]
\item
The set $\RR{3}$, which you can think of as ordinary space,
is the set of all ordered triples of real numbers:
\[
\R{3} = \br[\big]{(x,y, z): x, y, z \in \R{} }.
\]
\end{Itemize}
\vspace{-.12in}
\end{example}

To generalize $\R{2}$ and $\R{3}$ to higher dimensions, we first need
to discuss the concept of lists.

\DefinitionBox{list, length}
{\begin{Itemize}
\item
Suppose $n$ is a nonnegative integer.  A \emph{list}\index{list} of \emph{length}\index{length of list} $n$ is an ordered collection of $n$ elements (which might be numbers, other lists, or more abstract objects).
\item
Two lists are equal if and only if they have the same length and the same elements in the same order.
\end{Itemize}}

\mar[3]{Many
mathematicians call a list of length\/ $n$ an\/
$n$\marbf{-tuple}.\index{tuple}}

Lists are often written as elements separated by commas and surrounded by parentheses. Thus a list of length two is an ordered pair that might be written as $(a, b)$. A list of length three is an ordered triple that might be written as $(x,y,z)$. A list of length $n$ might look like this:
\[
(z_1, \dots, z_n).
\]

Sometimes we will use the word \emph{list} without specifying its length. Remember, however, that by definition each list has a finite length that is a nonnegative integer. Thus an object that looks like
$
(x_1, x_2, \dotsc )$,
which might be said to have infinite length, is not a list.

A list of length $0$ looks like this: $(\ )$.  We consider such an object to be
a list so that some of our theorems will not have trivial exceptions.

Lists differ from finite sets in two ways: in lists, order matters and  repetitions
have meaning; in sets, order and repetitions are irrelevant.

\begin{example}{lists versus sets}
\begin{Itemize}
\item
The lists $(3,5)$ and $(5,3)$ are not equal, but the sets $\{3,5\}$ and
$\{5,3\}$ are equal.
\item
The lists $(4,4)$ and $(4, 4, 4)$ are not equal
(they do not
have the same length), although the sets $\{4,4\}$ and $\{4, 4, 4\}$ both
equal the set $\{4\}$.
\end{Itemize}
\exampleitemend
\vspace{2bp}
\end{example}

    \subsection{$\F n$}

To define the higher-dimensional analogues of $\R{2}$ and $\RR{3}$, we
will simply replace $\R{}$ with $\F{}$ (which equals $\R{}$ or $\C{}$) and replace the $2$ or
$3$ with an arbitrary positive integer.

\NotationBox{$n$}{Fix a positive integer $n$ for the rest of this chapter.}

\DefinitionBox{$\F{n}$, coordinate}
{$\F{n}$\sindex{Fn@$\F{n}\1$} is the set of all
lists of length $n$ of elements of $\F{}$:
\[
\F{n} = \br[\big]{(x_1, \dots, x_n): x_k \in\F{} \mbox{\ for\ } k = 1, \dots, n }.
\]
For $(x_1, \dots, x_n) \in \F{n}$ and $k \in \{1, \dots, n\}$, we say
that $x_k$ is the $k^\nth$ \emph{coordinate}\index{coordinate} of $(x_1, \dots, x_n)$.}

If $\F{} = \R{}$ and $n$ equals $2$ or $3$, then the definition above of $\F{n}$ agrees with our previous notions of $\R{2}$
and $\RR{3}$.

\begin{example}{$\C{4}$}
$\C{4}$ is the set of all lists of four complex numbers:
\[
\C{4} = \br[\big]{ (z_1, z_2, z_3, z_4) : z_1, z_2, z_3, z_4 \in \C{} }.
\]
\exampledisplayend
\end{example}

\mar{Read\/ \marbf{Flatland: A Romance of Many Dimensions}, by Edwin A. Abbott, for an amusing account of how\/ $\R{3}$ would be perceived by creatures living in\/ $\RR{2}$. This novel, published in 1884, may help you imagine a
physical space  of four or more dimensions.\index{Abbott, Edwin A.}\index{Flatland@\textit{Flatland}}}

If $n \ge 4$, we cannot visualize $\R{n}$ as a physical object.  Similarly, $\C{1}$ can be thought of as a
plane, but for \mbox{$n \ge 2$}, the human brain cannot provide a full image of
$\CC{n}$. However, even if $n$ is large, we can perform algebraic manipulations
in $\F{n}$ as easily as in $\R{2}$
or $\RR{3}$.  For example, addition in $\F{n}$ is defined as follows.

\DefinitionBox[106]{addition in\/ $\F{n}$}
{\emph{Addition} in
$\F{n}$\index{addition!of vectors in $\F{n}$} is defined by adding corresponding coordinates:
\[
(x_1, \dots, x_n) + (y_1, \dots, y_n) = (x_1 + y_1, \dots, x_n + y_n).
\]}%

Often the mathematics of $\F{n}$ becomes cleaner if we use a
single letter to denote a list of $n$ numbers, without explicitly
writing the coordinates. For example, the next result is stated with $x$ and $y$ in $\F{n}$ even though the proof requires the more cumbersome notation of $(x_1, \dots, x_n)$ and $(y_1, \dots, y_n)$.

\ResultBox{commutativity of addition in $\F{n}$}
{If $x, y \in \FF{n}$, then $x + y = y + x$.\index{commutativity}}

\proof
Suppose $x = (x_1, \dots, x_n) \in \F n$ and $y = (y_1, \dots, y_n) \in \FF n$. Then
\begin{align*}
x + y &= (x_1, \dots, x_n) + (y_1, \dots, y_n )\\[6bp]
&= (x_1 + y_1, \dots, x_n + y_n) \\[6bp]
&= (y_1 + x_1, \dots, y_n + x_n) \\[6bp]
&= (y_1, \dots, y_n) + (x_1, \dots, x_n )\\[6bp]
&= y + x,
\end{align*}
where the second and fourth equalities above hold because of the definition of addition in $\F{n}$ and the third equality holds because of the usual commutativity of addition in $\F{}\3$.
\qed

\mar[3]{The symbol\/ \qedsymbol\ \,means ``end of proof\/''\!.\sindex{qed@\qedsymbol\hspace{2bp}}}

If a single letter is used to denote an element of $\FF{n}$, then the same letter with appropriate subscripts is often used when coordinates must be displayed. For example, if $x \in \FF{n}$, then letting $x$ equal $(x_1, \dots, x_n)$ is good notation, as shown in the proof above.  Even better, work with just~$x$ and avoid explicit coordinates when possible.

\NotationBox{$0$}
{\label{0vector}
Let $0$ denote the list of length $n$ whose coordinates are all $0$:
\[
0 = (0, \dots, 0).
\]}

Here we are using the symbol $0$ in two different ways---on the
left side of the equation above, the symbol $0$ denotes a list of length $n$, which is an element of $\FF n$, whereas on the
right side, each $0$ denotes a number.  This potentially confusing
practice actually causes no problems because the context should always make clear which $0$ is intended.

\begin{example}{context determines which\/ $0$ is intended}
Consider the statement that $0$ is an additive identity for~$\F{n}$:
\[
x + 0 = x \quad \text{for all } x \in \FF{n}.
\]
Here the $0$ above is the list defined in \ref{0vector}, not the number $0$, because
we have not defined the sum of an element of $\F{n}$ (namely, $x$) and the
number~$0$.
\end{example}

\pagebreak[3]

\marF[9]{Mathematica/Axes.pdf}{Elements of\/ $\R{2}$ can be thought of\\as points or as vectors.}

A picture can aid our intuition.  We will draw pictures in~$\R{2}$
because we can sketch this space on two-dimensional surfaces
such as paper and computer screens.  A typical element of $\R{2}$ is a
point $v = (a, b)$.
Sometimes we think of $v$ not as a point but
as an arrow starting at the origin and ending at $(a, b)$, as shown here.  When we think of an element of $\R{2}$ as an arrow, we refer to it as a \emph{vector}\index{vector}.

\marF{\includegraphics{Mathematica/Vector.pdf}}{A vector.}

When we think of vectors in $\R{2}$ as arrows, we can move an
arrow parallel to itself (not changing its length or direction) and still think of
it as the same vector. With that viewpoint, you will often gain better understanding by
dispensing with the coordinate axes and the explicit coordinates and just thinking of the vector, as shown in the figure here. The two arrows shown here have the same length and same direction, so we think of them as the same vector.

\mar{Mathematical  models of the economy can have thousands of
variables, say\/ $x_1, \dots,  x_{5000}$, which means that we must work
in\/ $\R{5000}$.  Such a  space cannot be dealt with geometrically. However, the
algebraic approach works well.  Thus our subject is called\/
\marbf{linear algebra}.}

Whenever we use pictures in $\R{2}$ or use the somewhat vague language of
points and vectors, remember that these are just aids to our understanding,
not substitutes for the actual mathematics that we will develop.  Although we
cannot draw good pictures in high-dimensional spaces, the elements of these
spaces are as rigorously defined as elements of~$\RR{2}$.

For example,
$\paren[\Big]{2, -3, 17, \pi, \sqrt{2}}$ is an element of $\RR{5}$, and we may casually refer to
it as a point in $\R{5}$ or a vector in $\R{5}$ without worrying about whether
the geometry of $\R{5}$ has any physical meaning.

Recall that we defined the sum of two elements of $\F{n}$ to be the element of
$\F{n}$ obtained by adding corresponding coordinates; see~\ref{106}.  As we will now see, addition has a simple geometric interpretation in the
special case of~$\RR{2}$.

\marF{\includegraphics{Mathematica/TriangleInequality2.pdf}}{The sum of two vectors.}

Suppose we have two vectors $u$ and $v$ in $\R{2}$ that we want to add.  Move the vector $v$ parallel to itself
so that its initial point coincides with the end point of the vector~$u$, as shown here.  The sum
$u+v$ then equals the vector whose initial point equals the initial point of $u$
and whose end point equals the end point of the vector $v$, as shown here.

In the next definition, the $0$ on the right side of the displayed equation is the list $0 \in \FF{n}$.

\DefinitionBox{additive inverse in $\F{n}$,\/ $-x$}
{For $x \in \FF{n}$, the \emph{additive inverse}\index{additive inverse!in $\F{n}$} of $x$, denoted by $-x$, is the vector $-x \in \F{n}$ such that
\[
x + (-x) = 0.
\]
Thus if $x= (x_1, \dots, x_n)$, then $-x = (-x_1, \dots, -x_n)$.}

\marF{\includegraphics{Mathematica/Minusx.pdf}}{A vector and its additive inverse.}

The additive inverse of a vector in $\R 2$ is the vector with the same length but pointing in the opposite direction. The figure here illustrates this way of thinking about the additive inverse in $\RR{2}$. As you can see, the vector labeled $-x$ has the same length as the vector labeled $x$ but points in the opposite direction.

Having dealt with addition in $\FF{n}$, we now turn to multiplication.  We could
\mbox{define} a multiplication in $\F{n}$ in a similar fashion, starting with two
elements of $\F{n}$ and getting another element of $\F{n}$ by multiplying
corresponding coordinates.  Experience shows that this definition is not useful
for our purposes.   Another type of multiplication, called scalar multiplication,
will be central to our subject.  Specifically, we need to define what it means to
multiply an element of $\F{n}$ by an element of $\F{}\3$.

\DefinitionBox{scalar multiplication in\/ $\F{n}$}
{The \emph{product}\index{product!of scalar and vector in $\F{n}$} of a number $\la$ and a vector in $\F{n}$ is computed by multiplying each coordinate of the vector by $\la$:\index{scalar multiplication}
\[
\la (x_1, \dots, x_n) = (\la x_1, \dots, \la x_n);
\]
here $\la \in \F{}$ and $(x_1, \dots, x_n) \in \FF{n}$.}

\mar{Scalar multiplication in\/ $\F n$ multiplies together a scalar and a
vector, getting a vector. In contrast, the dot product in\/
$\R{2}$ or\/ $\R{3}$ multiplies together two vectors and gets a scalar.
Generalizations of the dot product will become important in Chapter~\ref{89}.}

Scalar multiplication has a nice geometric interpretation in
$\RR{2}$. If $\la>0$ and $x \in \RR{2}$, then $\la x$ is
the vector that  points in the same direction as $x$ and whose length is $\la$
times the length of $x$.  In other words, to get $\la x$, we shrink or stretch $x$ by a factor of $\la$,
depending on whether $\la < 1$ or $\la > 1$.

\marF[4]{\includegraphics{Mathematica/PositiveMultiplication.pdf}}{Scalar multiplication.}

If $\la<0$ and $x \in \RR{2}$, then $\la x$ is the
vector that points in the direction opposite to that of $x$ and whose length is $|\la|$
times the  length of $x$, as shown here.

\vspace{.3in} % VisualFormatting

\subsection{Digression on Fields} \label{fields}

A \emph{field} is a set\index{field} containing at least two distinct elements called $0$ and $1$, along with operations of addition and multiplication satisfying all properties listed in \ref{complexarith}. Thus $\R{}$ and $\C{}$ are fields, as is the set of rational numbers along with the usual operations of addition and multiplication. Another example of a field is the set $\{0, 1\}$ with the usual operations of addition and multiplication except that $1 + 1$ is defined to equal $0$.

In this book we will not deal with fields other than $\R{}$ and $\C{}$. However, many of the definitions, theorems, and proofs in linear  algebra that work for the fields $\R{}$ and $\C{}$ also work without change for arbitrary fields. If you prefer to do so, throughout much of this book (except for Chapters \ref{89} and \ref{183}, which deal with inner product spaces) you can think of $\F{}$ as denoting an arbitrary field instead of $\R{}$ or~$\C{}$. For results (except in the inner product chapters) that have as a hypothesis that $\F{}$ is $\C{}$, you can probably replace that hypothesis with the hypothesis that $\F{}$ is an algebraically closed field, which means that every nonconstant polynomial with coefficients in $\F{}$ has a zero. A few results, such as Exercise \ref{1threesub} in Section \ref{1Subspaces}, require the hypothesis on $\F{}$ that $1 + 1 \ne 0$.

\begin{exercises}

\exercise
Show that $\al + \beta = \beta + \al$ for all $\al, \beta \in \C{}$.


\exercise
Show that $(\al + \beta) +  \la = \al + (\beta + \la)$ for all $\al, \beta, \la \in \C{}$.


\exercise
Show that $(\al \beta)\la = \al(\beta \la)$ for all $\al, \beta, \la \in \C{}$.


\exercise
Show that $\la(\al + \beta) = \la \al + \la \beta$ for all $\la, \al, \beta \in
\C{}$.


\exercise
Show that for every $\al \in \C{}$, there exists a unique $\beta \in \C{}$ such that
$\al + \beta = 0$.


\exercise
Show that for every $\al \in \C{}$ with $\al \ne 0$, there exists a unique $\beta \in \C{}$ such that $\al \beta = 1$.


\exercise
Show that
\[
\frac{-1 + \sqrt{3} \, i}{2}
\]
is a cube root of $1$ (meaning that its cube  equals~$1$).


\exercise
Find two distinct square roots of $i$.


\exercise
Find $x \in \R{4}$ such that
\[
(4, -3, 1, 7) + 2x = (5, 9, -6, 8).
\]
\exercisedisplayend


\exercise
Explain why there does not exist $\la \in \C{}$ such that
\[
\la (2 - 3i, 5 + 4i, -6 + 7i) = (12 - 5i, 7 + 22i, -32 - 9i).
\]
\exercisedisplayend


\exercise
Show that $(x + y) + z  = x + (y + z) $ for all $x, y, z \in \FF{n}$.


\exercise
Show that $(ab)x = a(bx)$ for all $x \in \F{n}$ and all $a, b \in \F{}\3$.


\exercise
Show that $1x = x$ for all $x \in \FF{n}$.


\exercise
Show that $\la(x + y) = \la x + \la y$ for all $\la \in \F{}$ and all $x, y \in \FF{n}$.


\exercise
Show that $(a + b)x = ax + bx$ for all $a, b \in \F{}$ and all $x \in \FF{n}$.


\end{exercises}

\QuoteAuthor{%
``Can you do addition?'' the White Queen asked. ``What's one and one and one and one and one and one and one and one and one and one?''\\
``I don't know,'' said Alice. ``I lost count.''}
{\textit{Through the Looking Glass}, Lewis Carroll\index{Carroll, Lewis}\index{Through the Looking@\textit{Through the Looking Glass}}}

\section{Definition of Vector Space} \label{1DefVS}

The motivation for the definition of a vector space comes from
properties of addition and scalar multiplication in~$\F{n}$:
Addition is commutative, associative, and has an
identity.  Every element has an additive inverse.  Scalar
multiplication is associative. Scalar multiplication by $1$ acts as
expected. Addition and scalar multiplication are connected by distributive  properties.

We will define a vector space to be a set $V$
with an addition and a scalar multiplication on $V$ that satisfy the properties
in the paragraph above.

\DefinitionBox{addition, scalar multiplication}{
\begin{Itemize}
\item
An \emph{addition}\index{addition!of vectors}
on a set $V$ is a function that assigns an
element~$u + v \in V$\\to each pair of elements $u, v \in V\1$.
\item
A \emph{scalar
multiplication}\index{scalar  multiplication}\index{product!of scalar and vector} on a set $V$ is a function that assigns
an element $\la v \in V$ to each $\la \in \F{}$ and each $v \in V\1$.
\end{Itemize}}

Now we are ready to give the formal definition of a vector space.

\vspace{-.05in}

\DefinitionBox[vectorspace]{vector space}{%
A \emph{vector space}\index{vector space} is a set $V$ along with an addition on
$V$ and a scalar multiplication on $V$ such that the following properties hold.

\vspace{-.1in}

\begin{description}\addtolength{\itemsep}{-3.5bp}%VisualFormatting
\item[commutativity] \ \\
$u + v = v + u$ for all $u, v \in V\1$.\index{commutativity}

\item[associativity] \ \\
$(u + v) + w = u + (v + w)$ and $(ab)v = a(bv)$ for all $u, v, w \in V$ and for all $a, b \in \F{}\3$.\index{associativity}

\item[additive identity] \ \\
There exists an element $0 \in V$ such that $v + 0 = v$ for all $v \in V\1$.

\item[additive inverse] \ \\
For every $v \in V\1$, there exists $w \in V$ such that $v + w = 0$.\index{additive inverse!in vector space}

\item[multiplicative identity] \ \\
$1v = v$ for all $v \in V\1$.

\item[distributive properties] \ \\
$a(u + v) = au + av$ and $(a + b)v = av + bv$ for all $a, b \in \F{}$ and all $u, v \in V\1$.\index{distributive property}
\end{description}
}

The following geometric language sometimes aids our intuition.

\DefinitionBox{vector, point}
{Elements of a vector space are called \emph{vectors}\index{vector} or
\emph{points}\index{point}.}

\pagebreak[3]

The scalar multiplication in a vector space
depends on $\F{}\3$.  Thus when we need to be precise, we will say that
$V$ is a \emph{vector space over\/ $\F{}$} instead of saying simply that
$V$ is a vector space.  For example, $\R{n}$ is a vector space over
$\R{}$, and $\C{n}$ is a vector space over~$\C{}$.

\DefinitionBox{real vector space, complex vector space}
{\begin{Itemize}
\item
A vector space over $\R{}$ is called a \emph{real vector space}\index{real vector space}.
\item
A vector space over $\C{}$ is called a \emph{complex vector space}\index{complex vector space}.
\end{Itemize}}

Usually the choice of $\F{}$ is either clear from the context or irrelevant. Thus we often assume that $\F{}$ is lurking in the background without specifically mentioning it.\looseness=-1

\mar{The simplest vector space is $\{0\}$, which contains only one point.}

With the usual operations of addition and scalar multiplication, $\F{n}$ is a vector space over $\F{}\3$, as you should verify.  The example of $\F{n}$ motivated our definition of vector space.

\begin{example}{$\F{\infty}$}
$\F{\infty}$\sindex{Finfty@$\F{\infty}\1$} is defined to be  the set of all sequences of elements of $\F{}$:
\[
\F{\infty} = \br[\big]{ (x_1, x_2, \dotsc) : x_k \in \F{} \mbox{\ for\ } k = 1, 2, \dotsc}.
\]
Addition and scalar multiplication on $\F{\infty}$ are defined as expected:
\begin{gather*}
(x_1, x_2, \dotsc) + (y_1, y_2, \dotsc) = (x_1 + y_1, x_2 + y_2, \dotsc),\\
\la(x_1, x_2, \dotsc) = (\la x_1, \la x_2, \dotsc).
\end{gather*}
With these definitions, $\F{\infty}$ becomes a vector space over $\F{}\3$, as you
should verify.  The additive identity in this vector space is the sequence of all $0$'s.
\end{example}

Our next example of a vector space involves a set of functions.

\NotationBox{$\F{S}$}
{\begin{Itemize}
\item
If $S$ is a set, then $\F{S}$ denotes the set of functions from $S$ to $\F{}\3$.
\item
For $f, g \in \FF{S}$, the \emph{sum} $f + g \in \F{S}$\index{addition!of functions} is the function defined by
\[
(f + g)(x) = f(x) + g (x)
\]
for all $x \in S$.
\item
For $\la \in \F{}$ and $f\in \FF{S}$, the \emph{product} $\la f \in \F{S}$ is the function defined by
\[
(\la f)(x) = \la f(x)
\]
for all $x \in S$.
\end{Itemize}}\sindex{FS@$\F{S}\1$}

\pagebreak[3]

As an example of the notation above, if $S$ is the interval $[0, 1]$ and $\F{} = \R{}$, then $\R{[0, 1]}$ is the set of real-valued functions on the interval $[0, 1]$.

You should verify all three bullet points in the next example.

\begin{example}{$\F{S}$ is a vector space}
\begin{Itemize}
\item
If $S$ is a nonempty set, then $\F{S}$ (with the operations of addition and scalar multiplication as defined above) is a vector space over $\F{}\3$.
\item
The additive identity of $\F{S}$ is the function $0: S \to \F{}$ defined by
\[
0(x) = 0
\]
for all $x \in S$.
\item
For $f \in \FF{S}$, the additive inverse of $f$ is the function $-f: S \to \F{}$ defined by
\[
(-f)(x) = -f(x)
\]
for all $x \in S$.
\end{Itemize}
\vspace{-.15in}
\end{example}

\mar{The elements of the vector space\/ $\R{[0,1]}$ are real-valued
functions on\/ $[0, 1]$, not lists.  In general, a vector space is an abstract
entity whose elements might be lists, functions, or weird objects.}

The vector space $\F{n}$ is a special case of the vector space $\F{S}$ because each $(x_1, \ldots, x_n) \in \F{n}$ can be thought of as a function $x$ from the set $\{1, 2, \dots, n\}$ to $\F{}$ by writing $x(k)$ instead of $x_k$ for the $k^\nth$ coordinate of $(x_1, \dots, x_n)$. In other words, we can think of $\F{n}$ as $\FF{\{1, 2, \dots, n\}}$. Similarly, we can think of $\F{\infty}$ as $\FF{\{1, 2, \dots \}}$.

Soon we will see further examples of vector spaces, but first we need to
develop some of the elementary properties of vector spaces.

The definition of a vector space requires it to have an additive identity. The next result states that this identity is unique.

\ResultBox{unique additive identity}
{A vector space has a unique additive identity.}

\proof
Suppose $0$ and $0'$ are both additive identities for some vector space~$V\1$.
Then
\[
0' = 0' + 0 = 0 + 0' = 0,
\]
where the first equality holds because $0$ is an additive identity, the second equality comes from commutativity, and the
third equality holds because $0'$ is an additive identity.  Thus $0' = 0$,
proving that $V$ has only one additive identity.
\qed

Each element $v$ in a vector space  has an additive inverse, an element
$w$ in the vector space such that $v + w = 0$.  The next result shows that
each element in a vector space has only one  additive inverse.

\ResultBox{unique additive inverse}
{Every element in a vector space has a unique additive inverse.}

\proof
Suppose $V$ is a vector space. Let $v \in V\1$.  Suppose $w$ and $w'$ are
additive  inverses of $v$.  Then
\[
w = w + 0 = w + (v + w') = (w + v) + w' = 0 + w' = w'\!.
\]
Thus $w = w'\!$, as desired.
\qed

Because additive inverses are unique, the following notation now makes sense.

\NotationBox{$-v$, $w - v$}
{
Let $v, w \in V\1$. Then\sindex{v1@$-v$}
\begin{Itemize}
\item
$-v$ denotes the additive inverse of $v$;\index{additive inverse!in vector space}
\item
$w - v$ is defined to be $w + (-v)$.
\end{Itemize}}

Almost all results in this book involve some vector space.  To avoid having to restate frequently that $V$ is a vector space, we now make the necessary declaration once and for all.

\NotationBox{$V$}{For the rest of this book, $V$ denotes a vector space over $\F{}\3$.\sindex{V@$V\1$}}

In the next result, $0$ denotes a scalar (the number $0 \in \F{}$) on the left side of the equation and a vector (the additive identity of $V$) on the right side of the equation.

\ResultBox[64]{the number $0$ times a vector}{$0v = 0$ for every\/ $v \in V\1$.}

\mar[9]{The result in \ref{64} involves the \mbox{additive} identity of $V$ and scalar multiplication. The only part of the  definition of a vector space that connects vector addition and scalar multiplication is the distributive property. Thus the distributive  property must be used in the proof of \ref{64}.\index{distributive property}}

\proof
For $v \in V\1$, we have
\[
0v = (0 + 0)v = 0v + 0v.
\]
Adding the additive inverse of $0v$ to both sides of the equation above gives
$0 = 0v$, as desired.
\qed

In the next result, $0$ denotes the additive identity of $V\1$.  Although  their proofs are similar, \ref{64} and \ref{65} are not identical.  More precisely, \ref{64} states that the product of the scalar $0$ and  any vector equals the vector~$0$, whereas \ref{65} states that the  product of any scalar and the vector $0$ equals the vector $0$.

\ResultBox[65]{a number times the vector\/ $0$}{$a0 = 0$ for every\/ $a \in \F{}\3$.}

\proof
For $a \in \F{}\3$, we have
\[
a0 = a(0 + 0) = a0 + a0.
\]
Adding the additive inverse of $a0$ to both sides of the equation above gives
$0 = a0$, as desired.
\qed

Now we show that if an element of $V$ is multiplied by the
scalar~$-1$, then the result is the
additive inverse of the element of~$V\1$.

\ResultBox[2]{the number\/ $-1$ times a vector}{%
$(-1)v = -v$ for every\/ $v \in V\1$.
}

\proof
For $v \in V\1$, we have
\[
v + (-1)v = 1v + (-1)v = \bigl(1 + (-1)\bigr)v = 0v = 0.
\]
This equation says that $(-1)v$, when added to $v$, gives $0$.  Thus
$(-1)v$ is the additive inverse of $v$, as desired.
\qed

\begin{exercises}

\exercise
Prove that $-(-v) = v$ for every $v \in V\1$.


\exercise
Suppose $a \in \F{}\3$, $v \in V\1$, and $av = 0$. Prove that $a = 0$ or $v = 0$.


\exercise
Suppose $v, w \in V\1$. Explain why there exists a unique $x \in V$ such that $v + 3x = w$.


\exercise
The empty set is not a vector space. The empty set fails to satisfy only one of the requirements listed in the definition of a vector space (\ref{vectorspace}). Which one?


\exercise
Show that in the definition of a vector space (\ref{vectorspace}), the additive inverse condition can be replaced with the condition that
\[
0v = 0 \text{ for all } v \in V\1.
\]
Here the $0$ on the left side is the number $0$, and the $0$ on the right side is the additive identity of $V\1$.
\exercisecomment{The phrase a ``condition can be replaced'' in a definition means that the collection of objects satisfying the definition is unchanged if the original condition is replaced with the new condition.}


\exercise
Let $\infty$ and $-\infty$ denote two distinct objects, neither of which is in $\R{}$. Define an addition and scalar multiplication on
$\R{} \cup \{\infty, -\infty\}$ as you could guess from the notation. Specifically, the sum and product of two real numbers is as usual, and for
$t \in \R{}$ define
\vspace{-.08in}
\[
t \infty =
\begin{cases}
-\infty & \text{if } t < 0,\\
0 & \text{if } t = 0,\\
\infty & \text{if } t > 0,
\end{cases}
\quad\quad
t (-\infty) =
\begin{cases}
\infty & \text{if } t < 0,\\
0 & \text{if } t = 0,\\
-\infty & \text{if } t > 0,
\end{cases}
\]
and
\begin{align*}
t + \infty &= \infty + t = \infty + \infty = \infty,\\[-1bp]
\quad t + (-\infty) &= (-\infty) + t  = (-\infty) + (-\infty) = -\infty,\\[-1bp]
\infty + (-\infty) &= (-\infty) + \infty = 0.
\end{align*}
With these operations of addition and scalar multiplication, is $\R{} \cup \{\infty,-\infty\}$ a vector space over $\R{}$? Explain.


\exercise
Suppose $S$ is a nonempty set. Let $V^S$ denote the set of functions from $S$ to~$V\1$. Define a natural addition and scalar multiplication on $V^S\!$, and show that $V^S$ is a vector space with these definitions.


\exercise
Suppose $V$ is a real vector space.
\vspace{-.075in}
\begin{Itemize}
\item
The \emph{complexification} of $V\1$, denoted by $V\1_{\C{}}$, equals $V\1 \times V\1$. An element of $V\1_{\C{}}$ is an ordered pair $(u, v)$, where $u, v \in V\1$, but we write this as $u + iv$.\index{complexification!of a vector space}\sindex{VC@$V\1_{\C{}}$}
\item
Addition on $V\1_{\C{}}$ is defined by\label{1VC}
\[
(u_1 + i v_1) + (u_2 + i v_2) = (u_1 + u_2) + i (v_1 + v_2)
\]
for all $u_1, v_1, u_2, v_2 \in V\1$.
\item
Complex scalar multiplication on~$V\1_{\C{}}$ is defined by
\[
(a + bi)(u + iv) = (au - bv) + i(av + bu)
\]
for all $a, b \in \R{}$ and all $u, v \in V\1$.
\end{Itemize}
Prove that with the definitions of addition and
scalar multiplication as above, $V\1_{\C{}}$ is a complex vector space.
\exercisecomment{Think of\/ $V$ as a subset of\/ $V\1_{\C{}}$ by identifying\/ $u \in V$ with\/
$u + i0$. The construction of\/ $V\1_{\C{}}$ from\/ $V$ can then be thought of as generalizing the construction of\/ $\C{n}$ from\/ $\RR{n}$.}


\end{exercises}

\section{Subspaces} \label{1Subspaces}

By considering subspaces, we can greatly expand our examples of vector spaces.

\DefinitionBox{subspace}{A subset $U$ of $V$ is called a \emph{subspace}\index{subspace}
of $V$ if $U$ is also a vector space with the same additive identity, addition, and scalar multiplication as on~$V\1$.}

\mar[3]{Some people  use the terminology\/ \mbox{\marbf{linear subspace}}\index{linear
subspace}, which means the  same as subspace.}

The next result gives the easiest way to check whether a subset of a vector space is a subspace.

\vspace{6bp}

\ResultBox[subspace]{conditions for a subspace}
{A subset $U$ of $V$ is a subspace of $V$ if and only if $U$  satisfies the following three conditions.
\begin{description}

\item[additive identity] \ \\
$0 \in U$.

\item[closed under addition]\index{closed under addition}  \ \\
$u, w \in U$ implies $u + w \in U$.

\item[closed under scalar multiplication]\index{closed
under scalar multiplication} \ \\
$a \in \F{}$ and $u \in U$ implies $au \in U$.
\end{description}}

\mar{The additive identity condition above could be replaced with the condition that\/ $U$ is nonempty \uppar{because then taking\/ $u \in U$ and multiplying it by\/ $0$ would imply that\/ $0 \in U\hspace{1bp}$}. However, if a subset\/ $U$ of\/ $V$ is indeed a subspace, then usually the quickest way to show that\/ $U$ is nonempty is to show that\/ $0 \in U$.}

\proof
If $U$ is a subspace of $V\1$, then $U$ satisfies the three conditions above by the definition of vector space.

Conversely, suppose $U$ satisfies the three conditions above. The first condition ensures that the additive identity of $V$ is in~$U$. The second condition ensures that addition makes
sense on~$U$. The third condition ensures that scalar multiplication makes sense
on~$U$.

If $u \in U$, then $-u$ \sbr{which equals $(-1)u$ by~\ref{2}} is also in $U$ by the third condition above. Hence
every element of $U$ has an additive inverse in $U$.

The other parts of the definition of a vector space, such as associativity and commutativity, are
automatically satisfied for $U$ because they hold on the
larger space~$V\1$. Thus $U$ is a vector space and hence is a subspace of $V\1$.
\qed

\pagebreak[3]

The three conditions in the result above usually enable us to determine quickly whether a
given subset of~$V$ is a subspace of~$V\1$. You should verify all assertions in the next example.

\begin{example}[subspaces]{subspaces}
\begin{subtheorem}*
\item
If $b \in \F{}\3$, then
\[
\br[\big]{(x_1, x_2, x_3, x_4) \in \F{4} : x_3 = 5 x_4 + b }
\]
is a subspace of $\F{4}$ if and only if $b = 0$.
\item
The set of continuous real-valued functions on the interval $[0, 1]$ is a subspace of $\RR{[0, 1]}$.
\item
The set of differentiable real-valued functions on $\R{}$ is a subspace of~$\RR{\R{}}$.
\item
The set of differentiable real-valued functions $f$ on the interval $(0, 3)$ such that $f'\!(2) = b$ is a subspace of~$\R{(0,3)}$ if and only if
$b = 0$.
\item
The set of all sequences of complex numbers with limit $0$ is a subspace of~$\CC{\infty}$.
\end{subtheorem}
\exampleitemend
\end{example}

\mar{The set\/ $\{0\}$ is the smallest subspace of\/ $V\1$,
and\/ $V$  itself is the largest subspace of\/ $V\1$. The empty set is not a
subspace of\/ $V$ because a subspace must be a vector space and hence must contain at least one element, namely, an additive \mbox{identity}.}

Verifying some of the items above shows the linear structure underlying parts of calculus. For example, (b) above requires the result that the sum of two continuous functions is continuous. As another example, (d) above requires the result that for a constant $c$, the derivative of $cf$ equals $c$ times the derivative of $f$.

The subspaces of $\R{2}$ are precisely $\{0\}$, all
lines in $\R{2}$ containing the origin, and $\RR{2}$.  The subspaces of $\R{3}$ are precisely $\{0\}$, all
lines in $\R{3}$ containing the origin, all planes in $\R{3}$ containing the origin, and $\RR{3}$.
To prove that all these objects are indeed subspaces is straightforward---the hard part is
to show that they are the only subspaces of $\R{2}$ and $\RR{3}$. That task will be
easier after we introduce some additional tools in the next chapter.

\subsection{Sums of Subspaces}

\mar[4]{The union of subspaces is rarely a subspace \uppar{see Exercise~\ref{218}}, which is why we usually work with sums rather than unions.}

When dealing with vector spaces, we are usually interested only in \mbox{subspaces}, as opposed to arbitrary
subsets. The notion of the sum of subspaces will be useful.

\smallskip %VisualFormatting

\DefinitionBox{sum of subspaces}
{Suppose $V\1_1, \dots, V\1_m$ are subspaces of~$V\1$.  The
\emph{sum} of $V\1_1, \dots, V\1_m$, denoted by $V\1_1 + \dots + V\1_m$, is the set of all possible sums of elements
of $V\1_1, \dots, V\1_m$.  More precisely,\index{addition!of subspaces}\index{sum of subspaces}
\[
V\1_1 + \dots + V\1_m = \{ v_1 + \dots + v_m: v_1 \in V\1_1, \dots, v_m \in V\1_m \}.
\]}

\pagebreak[3]

Let's look at some examples of sums of subspaces.

\begin{example}{a sum of subspaces of\/ $\F{3}$}
Suppose $U$ is the set of all elements of~$\F{3}$ whose second and third coordinates equal~$0$, and
$W$ is the set of all elements of~$\F{3}$ whose first and third coordinates equal~$0$:
\[
U = \br[\big]{(x, 0, 0) \in \F{3}: x \in \F{}} \andd W = \br[\big]{(0, y, 0) \in \F{3}: y \in \F{}}.
\]
Then
\[
U + W = \br[\big]{(x, y, 0) \in \F{3}: x, y \in \F{}},
\]
as you should verify.
\end{example}

\begin{example}{a sum of subspaces of\/ $\F{4}$}
Suppose
\[
U = \br[\big]{(x, x, y, y)\in\F{4}: x, y \in \F{}} \andd W = \br[\big]{(x, x, x, y) \in \F{4}: x, y \in \F{}}.
\]
Using words rather than symbols, we could say that $U$ is the set of elements of $\F4$ whose first two coordinates equal each other and whose third and fourth coordinates equal each other. Similarly, $W$ is the set of elements of $\F4$ whose first three coordinates equal each other.

To find a description of $U + W\3$, consider a typical element $(a, a, b, b)$ of $U$ and a typical element $(c,c, c, d)$ of $W\3$, where $a, b, c, d \in \F{}\3$. We have
\[
(a, a, b, b) + (c, c, c, d) = (a + c, a+c, b+c, b+d),
\]
which shows that every element of $U + W$ has its first two coordinates equal to each other. Thus
\begin{equation} \label{UpW}
U + W \subset \br[\big]{(x, x, y, z) \in \F{4}: x, y, z \in \F{}}.
\end{equation}

To prove the inclusion in the other direction, suppose $x, y, z \in \F{}\3$. Then
\[
(x, x, y, z) = (x,x,y,y) + (0, 0, 0, z-y),
\]
where the first vector on the right is in $U$ and the second vector on the right is in~$W\3$. Thus $(x, x, y, z) \in U + W\3$, showing that the inclusion \ref{UpW} also holds in the opposite direction. Hence
\[
U + W = \br[\big]{(x, x, y, z) \in \F{4}: x, y, z \in \F{}},
\]
which shows that $U + W$ is the set of elements of $\F4$ whose first two coordinates equal each other.
\end{example}

The next result states that the sum of subspaces is a subspace, and is in fact the smallest subspace containing all the summands (which means that every subspace containing all the summands also contains the sum).

\ResultBox{sum of subspaces is the smallest containing subspace}
{Suppose $V\1_1, \dots, V\1_m$ are subspaces of $V\1$. Then $V\1_1 + \dots + V\1_m$ is the smallest subspace
of~$V$ containing $V\1_1, \dots, V\1_m$.}

\proof
The reader can verify that $V\1_1 + \dots + V\1_m$ contains the additive identity $0$ and is closed under addition and scalar multiplication. Thus \ref{subspace} implies that $V\1_1 + \dots + V\1_m$ is a subspace of $V\1$.

\mar[10]{Sums of subspaces in the theory of vector spaces are analogous
to unions of subsets in set theory.  Given two subspaces of a vector space, the
smallest subspace containing them is their sum.  Analogously, given two
subsets of a set, the smallest subset containing them is their union.}

The subspaces $V\1_1, \dots, V\1_m$ are
all contained in $V\1_1 + \dots + V\1_m$ (to see this, consider sums
$v_1 + \dots + v_m$ where all except one of the $v_k$'s are~$0$).  Conversely, every
subspace of $V$ containing $V\1_1, \dots, V\1_m$ contains
$V\1_1 + \dots + V\1_m$ (because subspaces must contain all finite sums of their
elements). Thus $V\1_1 + \dots + V\1_m$ is the smallest subspace
of~$V$ containing $V\1_1, \dots, V\1_m$.
\qed

\subsection{Direct Sums}

Suppose $V\1_1, \dots, V\1_m$ are subspaces of $V\1$. Every element of $V\1_1 + \dots + V\1_m$ can be written in the
form
\[
v_1 + \dots + v_m,
\]
where each $v_k \in V\1_k$.  Of special interest are cases in which each
vector in $V\1_1 + \dots + V\1_m$ can be represented in the form above in only one way.  This situation is
so important that it gets a special name (direct sum) and a special symbol ($\oplus$).

\DefinitionBox[directsum]{direct sum,\/ $\oplus$\sindex{directsum@$\oplus$}}
{Suppose $V\1_1, \dots,  V\1_m$ are subspaces of $V\1$.
\begin{Itemize}
\item
The sum $V\1_1 + \dots + V\1_m$ is called a \emph{direct sum}\index{direct sum} if each element of~$V\1_1 + \dots + V\1_m$ can be written in only one way as a sum $v_1 + \dots + v_m$, where each $v_k \in V\1_k$.
\item
If $V\1_1 + \dots + V\1_m$ is a direct sum, then $V\1_1 \oplus \dots \oplus V\1_m$ denotes $V\1_1 + \dots + V\1_m$, with the $\oplus$ notation serving as an indication that this is a direct sum.
\end{Itemize}}

\begin{example}{a direct sum of two subspaces}
Suppose $U$ is the subspace
of~$\F{3}$ of those vectors whose last coordinate equals~$0$, and
$W$ is the subspace of~$\F{3}$ of those vectors whose first two
coordinates equal~$0$:
\[
U = \br[\big]{(x, y, 0) \in \F{3}: x, y \in \F{}} \andd W = \br[\big]{(0, 0, z) \in \F{3}: z \in \F{}}.
\]
Then $\F{3} = U \oplus W\3$, as you should verify.
\end{example}

\pagebreak[2]

\begin{example}{a direct sum of multiple subspaces}
\mar{To produce $\oplus$ in \TeX, type \textup{\texttt{\backslash oplus}}.}
Suppose $V\1_k$ is the subspace of $\F{n}$ of those vectors whose coordinates are all $0$, except possibly in the $k^\nth$ slot; for example, \mbox{$V\1_2 = \br[\big]{(0, x, 0, \dots, 0) \in \F{n}: x \in \F{}}$}.  Then
\[
\F{n} = V\1_1 \oplus \dots \oplus V\1_n,
\]
as you should verify.
\end{example}

Sometimes nonexamples add to our understanding as much as examples.

\begin{example}[3direct]{a sum that is not a direct sum}
Suppose
\begin{align*}
V\1_1 &= \br[\big]{(x, y, 0) \in \F{3}: x, y \in \F{}}, \\[6bp]
V\1_2 &= \br[\big]{(0, 0, z) \in \F{3}: z \in \F{}}, \\[6bp]
V\1_3 &= \br[\big]{(0, y, y) \in \F{3}: y \in \F{}}.
\end{align*}
Then $\F{3} = V\1_1 + V\1_2 + V\1_3$ because every vector
$(x, y, z) \in \F{3}$ can be written as
\[
(x, y, z) = (x, y, 0) + (0, 0, z) + (0, 0, 0),
\]
where the first vector on the right side is in~$V\1_1$, the second vector is
in~$V\1_2$, and the third vector is in~$V\1_3$.

However, $\F{3}$ does not equal the direct sum of $V\1_1, V\1_2 , V\1_3$, because the vector
$(0,0,0)$ can be written in more than one way as a sum $v_1 + v_2 + v_3$, with
each $v_k \in V\1_k$.   Specifically, we have
\[
(0,0,0) = (0,1,0) + (0,0,1) + (0,-1,-1)
\]
and, of course,
\[
(0,0,0) = (0,0,0) + (0,0,0) + (0,0,0),
\]
where the first vector on the right side of each equation above is in~$V\1_1$, the
second vector is in~$V\1_2$, and the third vector is in~$V\1_3$. Thus the sum $V\1_1 + V\1_2 + V\1_3$ is not a direct sum.
\end{example}

\mar[8]{The symbol\/ $\oplus$, which is a plus sign \mbox{inside} a circle, reminds us that we are dealing with a special
type of sum of subspaces---each element in the direct sum can be
represented in only one way as a sum of elements from the specified subspaces.}

The definition of direct sum requires every vector in the sum to have a unique representation as an appropriate sum. The next
result shows that when deciding whether a sum of subspaces is a direct sum, we only need to consider whether $0$ can be uniquely written as an appropriate sum.

\ResultBox[30]{condition for a direct sum}{%
Suppose $V\1_1, \dots, V\1_m$ are subspaces of $V\1$.
Then $V\1_1 + \dots + V\1_m$ is a direct sum if and  only if
the only way to write $0$ as a sum $v_1 + \dots + v_m$, where each
$v_k \in V\1_k$, is by taking each $v_k$ equal to\/ $0$.}

\proof
First suppose $V\1_1 + \dots + V\1_m$ is a direct sum.  Then the definition of direct sum implies that the only way to write $0$ as a sum
$v_1 + \dots + v_m$, where each $v_k \in V\1_k$, is by taking each $v_k$ equal to $0$.

Now suppose that the only way to write $0$ as a sum $v_1 + \dots + v_m$, where each
$v_k \in V\1_k$, is by taking each $v_k$ equal to $0$.  To show that $V\1_1 + \dots + V\1_m$ is a direct sum, let $v \in V\1_1 + \dots + V\1_m$.  We can write
\[
v = v_1 + \dots + v_m
\]
for some $v_1 \in V\1_1, \dots, v_m \in V\1_m$.  To show
that this representation is unique, suppose we also have
\[
v = u_1 + \dots + u_m,
\]
where $u_1 \in V\1_1, \dots, u_m \in V\1_m$.  Subtracting
these two equations, we have
\[
0 = (v_1 - u_1) + \dots + (v_m - u_m).
\]
Because $v_1 - u_1 \in V\1_1, \dots, v_m - u_m \in V\1_m$, the equation above implies that each $v_k - u_k$ equals $0$.  Thus $v_1 = u_1, \dots, v_m = u_m$, as desired.
\qed

\mar[3]{The symbol $\iff$\sindex{iff@$\iff$} used below means ``if and only if\/''\!; this symbol could also be read to mean ``is equivalent to''\!.}

The next result gives a simple condition for testing whether a sum of two subspaces is a direct sum.

\vspace{7bp}

\ResultBox[219]{direct sum of two subspaces}{%
Suppose $U$ and $W$ are subspaces of~$V\1$. Then
\[
U + W \text{ is a direct sum} \iff U \cap W = \{0\}.
\]}

\proof
First suppose that $U + W$ is a direct sum.  If $v \in U \cap W\3$, then $0 = v + (-v)$, where $v \in U$ and
$-v \in W\3$.  By the unique representation of~$0$ as the sum of a vector in~$U$
and a vector in~$W\3$, we have $v = 0$.  Thus $U \cap W = \{0\}$, completing
the proof in one direction.

To prove the other direction, now suppose $U \cap W =
\{0\}$.  To prove that $U + W$ is a direct sum, suppose $u \in U$, $w \in W\3$, and
\[
0 = u + w.
\]
To complete the proof, we only need to show that
$u = w = 0$ (by \ref{30}).   The equation above implies that $u = -w \in W\3$.  Thus
$u \in U \cap W\3$. Hence $u = 0$, which by the equation above implies
that $w = 0$, completing the proof.
\qed

\mar[9]{Sums of subspaces  are analogous to unions of subsets.
Similarly, direct sums of subspaces are analogous to disjoint
unions of subsets.  No two subspaces of a vector space can be
disjoint, because both contain\/ $0$.  So disjointness is replaced, at least
in the case of two subspaces, with the requirement that the intersection
equal\/ $\{0\}$.}

The result above deals only with the
case of two subspaces.  When
asking about a possible direct sum with more
than two subspaces, it is not enough to test that each pair of the subspaces
intersect only at~$0$.  To see this, consider Example~\ref{3direct}.  In that nonexample of a direct sum, we have
$V\1_1 \cap V\1_2 = V\1_1 \cap V\1_3 = V\1_2 \cap V\1_3 = \{0\}$.

\begin{exercises}

\exercise
For each of the following subsets of $\FF{3}$, determine whether it is a
subspace of $\F{3}$.
\begin{subtheorem}*
\item
$\br[\big]{ (x_1, x_2, x_3) \in \F{3}: x_1 + 2 x_2 + 3 x_3 = 0 }$

\item
$\br[\big]{ (x_1, x_2, x_3) \in \F{3}: x_1 + 2 x_2 + 3 x_3 = 4 }$

\item
$\br[\big]{ (x_1, x_2, x_3) \in \F{3}: x_1 x_2 x_3 = 0 }$

\item
$\br[\big]{ (x_1, x_2, x_3) \in \F{3}: x_1 = 5 x_3 }$
\end{subtheorem}
\exercisesubtheoremend


\exercise
Verify all assertions about subspaces in Example \ref{subspaces}.


\exercise
Show that the set of differentiable real-valued functions $f$ on the interval $(-4, 4)$ such that $f'\!(-1) = 3f(2)$ is a subspace of
$\RR{(-4, 4)}$.


\pagebreak[1]

\exercise
Suppose $b \in \R{}$. Show that the set of continuous real-valued functions $f$ on the interval $[0, 1]$ such that $\int_0^1 f = b$ is a subspace of $\R{[0, 1]}$ if and only if $b = 0$.


\exercise
Is $\R{2}$ a subspace of the complex vector space $\C{2}$?


\exercise
\begin{SUBtheorem}
\item
Is $\br[\big]{(a, b, c) \in \R{3}: a^3 =b^3}$ a subspace of $\R{3}$?
\item
Is $\br[\big]{(a, b, c) \in \C{3}: a^3 =b^3}$ a subspace of $\C{3}$?
\end{SUBtheorem}


\exercise
Prove or give a counterexample: If $U$ is a nonempty subset of $\R{2}$ such that $U$ is closed
under addition and under taking additive inverses (meaning $-u \in U$
whenever $u \in U$), then $U$ is a subspace of~$\RR{2}$.


\exercise
Give an example of a nonempty subset $U$ of $\R{2}$ such that $U$ is closed
under scalar multiplication, but $U$ is not a subspace of~$\RR{2}$.


\exercise
A function $\fun f {\R{}} {\R{}}$ is called \emph{periodic} if there exists a positive number $p$ such that $f(x) = f(x + p)$ for all $x \in \R{}$. Is the set of periodic functions from $\R{}$ to $\R{}$ a subspace of $\R{\R{}}$? Explain.


\exercise
Suppose $V\1_1$ and $V\1_2$ are subspaces of $V\1$. Prove that the intersection $V\1_1 \cap V\1_2$ is a subspace of $V\1$.


\exercise
Prove that the intersection of every collection of subspaces of $V$ is a subspace of $V\1$.


\exercise
Prove that the union of two subspaces of $V$ is a subspace of $V$ if and only if
one of the subspaces is contained in the other. \label{218}


\exercise
Prove that the union of three subspaces of $V$ is a subspace of $V$ if and only if one of the subspaces contains the other two.\label{1threesub}
\exercisecomment{This exercise is surprisingly harder than Exercise \ref{218}, possibly because this exercise is not true if we replace $\F{}$ with a field containing only two elements.}


\exercise
Suppose
\[
U = \br[\big]{ (x, -x, 2x) \in \F3 : x\in \F{}} \andd W = \br[\big]{ (x, x, 2x) \in \F3: x \in \F{}}.
\]
Describe $U+W$ using symbols, and also give a description of $U + W$ that uses no symbols.


\exercise
Suppose $U$ is a subspace of $V\1$.  What is $U + U$?


\exercise
Is the operation of addition on the subspaces of $V$ commutative?
In other words, if $U$ and $W$ are subspaces of $V\1$, is
$U + W = W + U$?\index{commutativity}


\exercise
Is the operation of addition on the subspaces of $V$ associative?
In other words, if $V\1_1,V\1_2, V\1_3$ are subspaces of $V\1$, is
\[
(V\1_1 + V\1_2) + V\1_3 = V\1_1 + (V\1_2 + V\1_3)?
\]
\exercisedisplayend


\exercise
Does the operation of addition on the subspaces of $V$ have an
additive identity?  Which subspaces have additive inverses?


\exercise
Prove or give a counterexample: If $V\1_1, V\1_2, U$ are subspaces of~$V$ such
that
\[
V\1_1 + U = V\1_2 + U,
\]
then $V\1_1 = V\1_2$.


\exercise
Suppose
\[
U = \br[\big]{(x, x, y, y) \in \F{4}: x, y \in \F{}}.
\]
Find a subspace $W$ of $\F{4}$ such that
$\F{4} = U \oplus W\3$.


\exercise
Suppose
\[
U = \br[\big]{(x, y, x+y, x-y, 2x) \in \F{5}: x, y \in \F{}}.
\]
Find a subspace $W$ of $\F{5}$ such that
$\F{5} = U \oplus W\3$.


\exercise
Suppose
\[
U = \br[\big]{(x, y, x+y, x-y, 2x) \in \F{5}: x, y \in \F{}}.
\]
Find three subspaces $W\3_1, W\3_2, W\3_3$ of $\FF{5}$, none of which equals $\{0\}$, such that
$\F{5} = U \oplus W\3_1 \oplus W\3_2 \oplus W\3_3$.


\exercise
Prove or give a counterexample: If $V\1_1, V\1_2, U$ are subspaces of~$V$ such that
\[
V = V\1_1 \oplus U \andd V = V\1_2 \oplus U,
\]
then $V\1_1 = V\1_2$.
\hint{When trying to discover whether a conjecture in linear algebra is true or false, it is often useful to start by experimenting in\/ $\FF2$.}


\exercise
A function $f\colon \R{} \to \R{}$ is called \emph{even} if
\[
f(-x) = f(x)
\]
for all $x \in \R{}$. A function $f\colon \R{} \to \R{}$ is called \emph{odd} if
\[
f(-x) = -f(x)
\]
for all $x \in \R{}$. Let $V\1_\text{e}$ denote the set of real-valued even functions on $\R{}$ and let $V\1_\text{o}$ denote the set of real-valued odd functions on $\R{}$.
Show that $\R{\R{}} = V\1_\text{e} \oplus V\1_\text{o}$.


\end{exercises}

\CrCom

\endinput 
